\chapter{Effects}\label{ch:effects}

A sound carries three effect blocks --- \textsc{effect1}, \textsc{effect2} and
\textsc{reverb} --- and each is a \textsc{type} byte followed by twenty \textsc{value}
bytes, an \textsc{another eff send} and a \textsc{dynamic control}. The \textsc{type} byte
selects one of 56 effects, and what the twenty \textsc{value} bytes mean depends entirely
on which.

\section{What the type byte selects}
The numbers are not contiguous: they are grouped, with modulation effects below
\bytes{10}, reverbs at \bytes{10}--\bytes{1B}, distortion and dynamics at
\bytes{20}--\bytes{27}, a second modulation group at \bytes{30}--\bytes{38}, two-effect
combinations at \bytes{40}--\bytes{4B} and three-effect combinations at
\bytes{60}--\bytes{63}.

Twelve of the reverbs are marked \textsc{(rev only)} in the published guide: they are
intended for the \textsc{reverb} block rather than \textsc{effect1} or \textsc{effect2}.

\begin{note}{Two sources agree on which effects exist}
The table below was transcribed from the guide. Independently, this project
disassembled the effects DSP and produced a manifest of every distinct effect program
with the number that selects it. The two sets of numbers are equal --- 56 effects,
no more on either side --- which is what makes the numbering here trustworthy.

They disagree only in spelling, because the ROM carries the short names the instrument
displays and the guide prints long ones. The table gives the guide's.
\end{note}

\section{What the value bytes mean}
The \textsc{value} list below is the order the bytes appear in, so the first name is
\textsc{value1}, the second \textsc{value2}, and so on. An effect that names fewer than
twenty leaves the rest unused.

The combination effects concatenate their components: the first block's parameters, then
a dry/wet for the second block, then the second block's. A three-effect combination
continues the same way.

\begin{caution}{The value lists are transcribed, and not checked by anything}
The effect \emph{numbers} are confirmed twice over. The \emph{names} in the value lists
are not: they were read from scans, nothing independent states them, and a misreading
would survive. Treat them as a guide to what the bytes are for rather than as a
specification.

Two further things the guide gives and this edition does not reproduce: the conversion
from a byte value to a physical unit --- seconds, hertz, decibels --- which it prints as
eighteen tables, and the block diagram of each effect.
\end{caution}

%% GENERATED by notes/sysex-probes/sysex_effects.py --tex
%% Source: effects.json, effect_value_numbers.json and
%% dsp/programs.tsv. Do not edit.

{\small
\begin{longtable}{ll>{\raggedright\arraybackslash}p{40mm}>{\raggedright\arraybackslash}p{78mm}}
\toprule
no. & n & effect & parameters, with the \textsc{value} byte each occupies \\
\midrule\endfirsthead
\toprule no. & n & effect & parameters, with the \textsc{value} byte each occupies \\
\midrule\endhead
\bytes{00} & 0 & NO OPERATION & --- \\
\bytes{01} & 5 & CHORUS & WET~\textbf{1}, DEPTH~\textbf{2}, LFO SPEED~\textbf{3}, LFO WAVEFORM~\textbf{4}, VOLUME~\textbf{5} \\
\bytes{02} & 7 & MODULATED CHORUS & WET~\textbf{1}, DEPTH~\textbf{2}, SLOW LFO SPEED~\textbf{3}, FAST LFO SPEED~\textbf{4}, FAST LFO BALANCE~\textbf{5}, LFO WAVEFORM~\textbf{6}, VOLUME~\textbf{7} \\
\bytes{03} & 9 & ENHANCER & WET~\textbf{1}, MANUAL~\textbf{2}, LOW MIX~\textbf{3}, HIGH MIX~\textbf{4}, DELAY TIME L~\textbf{5--6}, DELAY TIME R~\textbf{7--8}, VOLUME~\textbf{9} \\
\bytes{04} & 8 & FLANGER & WET~\textbf{1}, DEPTH~\textbf{2}, LFO SPEED~\textbf{3}, RESONANCE~\textbf{4}, MANUAL~\textbf{5}, PHASE~\textbf{6}, LFO WAVEFORM~\textbf{7}, VOLUME~\textbf{8} \\
\bytes{05} & 8 & PHASER & WET~\textbf{1}, DEPTH~\textbf{2}, LFO SPEED~\textbf{3}, RESONANCE~\textbf{4}, MANUAL~\textbf{5}, PHASE~\textbf{6}, LFO WAVEFORM~\textbf{7}, VOLUME~\textbf{8} \\
\bytes{06} & 5 & ENSEMBLE & WET~\textbf{1}, DEPTH~\textbf{2}, LFO SPEED~\textbf{3}, LFO WAVEFORM~\textbf{4}, VOLUME~\textbf{5} \\
\bytes{08} & 6 & GATED REVERB & WET~\textbf{1}, GATE TIME~\textbf{2}, HIGH DUMP GAIN~\textbf{3}, THRESHOLD~\textbf{4}, MASK TIME~\textbf{5}, VOLUME~\textbf{6} \\
\bytes{09} & 9 & SINGLE DELAY & WET~\textbf{1}, DELAY L~\textbf{2--3}, DELAY R~\textbf{4--5}, FEEDBACK L~\textbf{6}, FEEDBACK R~\textbf{7}, HIGH DUMP GAIN~\textbf{8}, VOLUME~\textbf{9} \\
\bytes{0A} & 16 & MULTI TAP DELAY & WET~\textbf{1}, DELAY 1~\textbf{2--3}, DELAY 2~\textbf{4--5}, DELAY 3~\textbf{6--7}, DELAY 4~\textbf{8--9}, PAN 1~\textbf{10}, PAN 2~\textbf{11}, PAN 3~\textbf{12}, PAN 4~\textbf{13}, FEED BACK~\textbf{14}, HIGH DUMP GAIN~\textbf{15}, VOLUME~\textbf{16} \\
\bytes{0B} & 10 & MANUAL DELAY & WET~\textbf{1}, MODULATION DEPTH~\textbf{2}, DELAY L~\textbf{3--4}, DELAY R~\textbf{5--6}, FEEDBACK L~\textbf{7}, FEEDBACK R~\textbf{8}, HIGH DUMP GAIN~\textbf{9}, VOLUME~\textbf{10} \\
\bytes{10} & 5 & ROOM REVERB 1 (REV only) & REVERB TIME~\textbf{1}, PRE DELAY~\textbf{2}, HIGH DUMP GAIN~\textbf{3}, EARLY REFL LEVEL~\textbf{4}, VOLUME~\textbf{5} \\
\bytes{11} & 5 & ROOM REVERB 2 (REV only) & REVERB TIME~\textbf{1}, PRE DELAY~\textbf{2}, HIGH DUMP GAIN~\textbf{3}, EARLY REFL LEVEL~\textbf{4}, VOLUME~\textbf{5} \\
\bytes{12} & 5 & PLATE REVERB 1 (REV only) & REVERB TIME~\textbf{1}, PRE DELAY~\textbf{2}, HIGH DUMP GAIN~\textbf{3}, EARLY REFL LEVEL~\textbf{4}, VOLUME~\textbf{5} \\
\bytes{13} & 5 & PLATE REVERB 2 (REV only) & REVERB TIME~\textbf{1}, PRE DELAY~\textbf{2}, HIGH DUMP GAIN~\textbf{3}, EARLY REFL LEVEL~\textbf{4}, VOLUME~\textbf{5} \\
\bytes{14} & 5 & CONCERT REVERB 1 (REV only) & REVERB TIME~\textbf{1}, PRE DELAY~\textbf{2}, HIGH DUMP GAIN~\textbf{3}, EARLY REFL LEVEL~\textbf{4}, VOLUME~\textbf{5} \\
\bytes{15} & 5 & CONCERT REVERB 2 (REV only) & REVERB TIME~\textbf{1}, PRE DELAY~\textbf{2}, HIGH DUMP GAIN~\textbf{3}, EARLY REFL LEVEL~\textbf{4}, VOLUME~\textbf{5} \\
\bytes{16} & 5 & DARK REVERB 1 (REV only) & REVERB TIME~\textbf{1}, PRE DELAY~\textbf{2}, HIGH DUMP GAIN~\textbf{3}, EARLY REFL LEVEL~\textbf{4}, VOLUME~\textbf{5} \\
\bytes{17} & 5 & DARK REVERB 2 (REV only) & REVERB TIME~\textbf{1}, PRE DELAY~\textbf{2}, HIGH DUMP GAIN~\textbf{3}, EARLY REFL LEVEL~\textbf{4}, VOLUME~\textbf{5} \\
\bytes{18} & 5 & BRIGHT REVERB 1 (REV only) & REVERB TIME~\textbf{1}, PRE DELAY~\textbf{2}, HIGH DUMP GAIN~\textbf{3}, EARLY REFL LEVEL~\textbf{4}, VOLUME~\textbf{5} \\
\bytes{19} & 5 & BRIGHT REVERB 2 (REV only) & REVERB TIME~\textbf{1}, PRE DELAY~\textbf{2}, HIGH DUMP GAIN~\textbf{3}, EARLY REFL LEVEL~\textbf{4}, VOLUME~\textbf{5} \\
\bytes{1A} & 5 & WAVE REVERB 1 (REV only) & REVERB TIME~\textbf{1}, PRE DELAY~\textbf{2}, HIGH DUMP GAIN~\textbf{3}, EARLY REFL LEVEL~\textbf{4}, VOLUME~\textbf{5} \\
\bytes{1B} & 5 & WAVE REVERB 2 (REV only) & REVERB TIME~\textbf{1}, PRE DELAY~\textbf{2}, HIGH DUMP GAIN~\textbf{3}, EARLY REFL LEVEL~\textbf{4}, VOLUME~\textbf{5} \\
\bytes{20} & 4 & DISTORTION & WET~\textbf{1}, DRIVE~\textbf{2}, ADJUST~\textbf{3}, VOLUME~\textbf{4} \\
\bytes{21} & 4 & OVERDRIVE & WET~\textbf{1}, DRIVE~\textbf{2}, ADJUST~\textbf{3}, VOLUME~\textbf{4} \\
\bytes{22} & 4 & FUZZ & WET~\textbf{1}, DRIVE~\textbf{2}, ADJUST~\textbf{3}, VOLUME~\textbf{4} \\
\bytes{23} & 7 & EXCITER & WET~\textbf{1}, DRIVE~\textbf{2}, ADJUST~\textbf{3}, EMPHASIS Fc~\textbf{4--5}, EMPHASIS GAIN~\textbf{6}, VOLUME~\textbf{7} \\
\bytes{24} & 6 & COMPRESSOR & WET~\textbf{1}, THRESHOLD~\textbf{2}, RATIO~\textbf{3}, ATTACK SENSITIVITY~\textbf{4}, RELEASE SENSITIVITY~\textbf{5}, VOLUME~\textbf{6} \\
\bytes{25} & 5 & SLOW ATTACKER & WET~\textbf{1}, THRESHOLD~\textbf{2}, ATTACK RATE~\textbf{3}, RELEASE RATE~\textbf{4}, VOLUME~\textbf{5} \\
\bytes{26} & 1 & NOISE GENERATOR & VOLUME~\textbf{1} \\
\bytes{27} & 13 & PARAMETRIC EQ & BAND EMPHASIS 1 Fc~\textbf{1--2}, BAND EMPHASIS 1 Q~\textbf{1--2}, BAND EMPHASIS 1 G~\textbf{1--2}, BAND EMPHASIS 2 Fc~\textbf{3--4}, BAND EMPHASIS 2 Q~\textbf{3--4}, BAND EMPHASIS 2 G~\textbf{3--4}, BAND EMPHASIS 3 Fc~\textbf{5--6}, BAND EMPHASIS 3 Q~\textbf{5--6}, BAND EMPHASIS 3 G~\textbf{5--6}, BAND EMPHASIS 4 Fc~\textbf{7--8}, BAND EMPHASIS 4 Q~\textbf{7--8}, BAND EMPHASIS 4 G~\textbf{7--8}, BAND EMPHASIS 5 Fc~\textbf{9--10}, BAND EMPHASIS 5 Q~\textbf{9--10}, BAND EMPHASIS 5 G~\textbf{9--10}, BAND EMPHASIS 6 Fc~\textbf{11--12}, BAND EMPHASIS 6 Q~\textbf{11--12}, BAND EMPHASIS 6 G~\textbf{11--12}, VOLUME~\textbf{13} \\
\bytes{30} & 6 & AUTO PAN & WET~\textbf{1}, DEPTH~\textbf{2}, LFO SPEED~\textbf{3}, PHASE~\textbf{4}, LFO WAVEFORM~\textbf{5}, VOLUME~\textbf{6} \\
\bytes{31} & 6 & PITCH SHIFTER & WET~\textbf{1}, PITCH L~\textbf{2}, PITCH R~\textbf{3}, PRE DELAY~\textbf{4}, FEEDBACK~\textbf{5}, VOLUME~\textbf{6} \\
\bytes{32} & 6 & VIBRATO & WET~\textbf{1}, DEPTH~\textbf{2}, LFO SPEED~\textbf{3}, PHASE~\textbf{4}, LFO WAVEFORM~\textbf{5}, VOLUME~\textbf{6} \\
\bytes{33} & 6 & PEDAL WAH & WET~\textbf{1}, RESONANCE~\textbf{2}, MANUAL~\textbf{3}, SWEEP RANGE~\textbf{4}, WAH CENTER Fc~\textbf{5}, VOLUME~\textbf{6} \\
\bytes{34} & 5 & AUTO WAH & WET~\textbf{1}, RESONANCE~\textbf{2}, MANUAL~\textbf{3}, SWEEP RANGE~\textbf{4}, VOLUME~\textbf{5} \\
\bytes{35} & 15 & ROTARY SPEAKER & WET~\textbf{1}, DRIVE~\textbf{2}, VOLUME ADJUST~\textbf{3}, TREBLE DEPTH~\textbf{4}, TREBLE FAST~\textbf{5}, TREBLE SLOW~\textbf{6}, TREBLE WIND UP~\textbf{7}, TREBLE WIND DOWN~\textbf{8}, BASS DEPTH~\textbf{9}, BASS FAST~\textbf{10}, BASS SLOW~\textbf{11}, BASS WIND UP~\textbf{12}, BASS WIND DOWN~\textbf{13}, VOLUME~\textbf{14}, SLOW/FAST~\textbf{15} \\
\bytes{36} & 5 & RING MODULATOR & WET~\textbf{1}, OSC SPEED~\textbf{2}, PHASE~\textbf{3}, OSC WAVEFORM~\textbf{4}, VOLUME~\textbf{5} \\
\bytes{37} & 8 & HAAS EFFECT & WET~\textbf{1}, DELAY TIME L~\textbf{2--3}, DELAY TIME R~\textbf{4--5}, BALANCE L~\textbf{6}, BALANCE R~\textbf{7}, VOLUME~\textbf{8} \\
\bytes{38} & 8 & MIX UP & WET~\textbf{1}, DEPTH~\textbf{2}, SLOW LFO SPEED~\textbf{3}, FAST LFO SPEED L~\textbf{4}, FAST LFO SPEED R~\textbf{5}, PHASE~\textbf{6}, LFO WAVEFORM~\textbf{7}, VOLUME~\textbf{8} \\
\bytes{40} & 12 & SINGLE DELAY + CHORUS & DELAY WET~\textbf{1}, DELAY L~\textbf{2--3}, DELAY R~\textbf{4--5}, FEEDBACK L~\textbf{6}, FEEDBACK R~\textbf{7}, CHORUS DRY/WET~\textbf{8}, DEPTH~\textbf{9}, LFO SPEED~\textbf{10}, LFO WAVEFORM~\textbf{11}, VOLUME~\textbf{12} \\
\bytes{41} & 11 & SINGLE DELAY + SINGLE DELAY & DELAY 1 WET~\textbf{1}, DELAY L~\textbf{2}, DELAY R~\textbf{3}, FEEDBACK L~\textbf{4}, FEEDBACK R~\textbf{5}, DELAY 2 DRY/WET~\textbf{6}, DELAY L~\textbf{7}, DELAY R~\textbf{8}, FEEDBACK L~\textbf{9}, FEEDBACK R~\textbf{10}, VOLUME~\textbf{11} \\
\bytes{42} & 15 & SINGLE DELAY + FLANGER & DELAY WET~\textbf{1}, DELAY L~\textbf{2--3}, DELAY R~\textbf{4--5}, FEEDBACK L~\textbf{6}, FEEDBACK R~\textbf{7}, FLANGER DRY/WET~\textbf{8}, DEPTH~\textbf{9}, LFO SPEED~\textbf{10}, RESONANCE~\textbf{11}, MANUAL~\textbf{12}, PHASE~\textbf{13}, LFO WAVEFORM~\textbf{14}, VOLUME~\textbf{15} \\
\bytes{43} & 13 & SINGLE DELAY + VIBRATO & DELAY WET~\textbf{1}, DELAY L~\textbf{2--3}, DELAY R~\textbf{4--5}, FEEDBACK L~\textbf{6}, FEEDBACK R~\textbf{7}, VIBRATO DRY/WET~\textbf{8}, DEPTH~\textbf{9}, LFO SPEED~\textbf{10}, PHASE~\textbf{11}, LFO WAVEFORM~\textbf{12}, VOLUME~\textbf{13} \\
\bytes{44} & 15 & SINGLE DELAY + PHASER & DELAY WET~\textbf{1}, DELAY L~\textbf{2--3}, DELAY R~\textbf{4--5}, FEEDBACK L~\textbf{6}, FEEDBACK R~\textbf{7}, PHASER DRY/WET~\textbf{8}, DEPTH~\textbf{9}, LFO SPEED~\textbf{10}, RESONANCE~\textbf{11}, MANUAL~\textbf{12}, PHASE~\textbf{13}, LFO WAVEFORM~\textbf{14}, VOLUME~\textbf{15} \\
\bytes{45} & 13 & PEDAL WAH + DELAY & WAH WET~\textbf{1}, RESONANCE~\textbf{2}, MANUAL~\textbf{3}, SWEEP RANGE~\textbf{4}, WAH CENTER Fc~\textbf{5}, DELAY DRY/WET~\textbf{6}, DELAY L~\textbf{7--8}, DELAY R~\textbf{9--10}, FEEDBACK L~\textbf{11}, FEEDBACK R~\textbf{12}, VOLUME~\textbf{13} \\
\bytes{46} & 12 & AUTO WAH + SINGLE DELAY & WAH WET~\textbf{1}, RESONANCE~\textbf{2}, MANUAL~\textbf{3}, SWEEP RANGE~\textbf{4}, DELAY DRY/WET~\textbf{5}, DELAY L~\textbf{6--7}, DELAY R~\textbf{8--9}, FEEDBACK L~\textbf{10}, FEEDBACK R~\textbf{11}, VOLUME~\textbf{12} \\
\bytes{47} & 9 & PEQ + CHORUS & BAND EMPHASIS 1 Fc~\textbf{1--2}, BAND EMPHASIS 1 Q~\textbf{1--2}, BAND EMPHASIS 1 G~\textbf{1--2}, BAND EMPHASIS 2 Fc~\textbf{3--4}, BAND EMPHASIS 2 Q~\textbf{3--4}, BAND EMPHASIS 2 G~\textbf{3--4}, CHORUS DRY/WET~\textbf{5}, DEPTH~\textbf{6}, LFO SPEED~\textbf{7}, LFO WAVEFORM~\textbf{8}, VOLUME~\textbf{9} \\
\bytes{48} & 12 & PEQ + SINGLE DELAY & BAND EMPHASIS 1 Fc~\textbf{1--2}, BAND EMPHASIS 1 Q~\textbf{1--2}, BAND EMPHASIS 1 G~\textbf{1--2}, BAND EMPHASIS 2 Fc~\textbf{3--4}, BAND EMPHASIS 2 Q~\textbf{3--4}, BAND EMPHASIS 2 G~\textbf{3--4}, DELAY DRY/WET~\textbf{5}, DELAY L~\textbf{6--7}, DELAY R~\textbf{8--9}, FEEDBACK L~\textbf{10}, FEEDBACK R~\textbf{11}, VOLUME~\textbf{12} \\
\bytes{49} & 12 & PEQ + FLANGER & BAND EMPHASIS 1 Fc~\textbf{1--2}, BAND EMPHASIS 1 Q~\textbf{1--2}, BAND EMPHASIS 1 G~\textbf{1--2}, BAND EMPHASIS 2 Fc~\textbf{3--4}, BAND EMPHASIS 2 Q~\textbf{3--4}, BAND EMPHASIS 2 G~\textbf{3--4}, FLANGER DRY/WET~\textbf{5}, DEPTH~\textbf{6}, LFO SPEED~\textbf{7}, RESONANCE~\textbf{8}, MANUAL~\textbf{9}, PHASE~\textbf{10}, LFO WAVEFORM~\textbf{11}, VOLUME~\textbf{12} \\
\bytes{4A} & 10 & PEQ + VIBRATO & BAND EMPHASIS 1 Fc~\textbf{1--2}, BAND EMPHASIS 1 Q~\textbf{1--2}, BAND EMPHASIS 1 G~\textbf{1--2}, BAND EMPHASIS 2 Fc~\textbf{3--4}, BAND EMPHASIS 2 Q~\textbf{3--4}, BAND EMPHASIS 2 G~\textbf{3--4}, VIBRATO DRY/WET~\textbf{5}, DEPTH~\textbf{6}, LFO SPEED~\textbf{7}, PHASE~\textbf{8}, LFO WAVEFORM~\textbf{9}, VOLUME~\textbf{10} \\
\bytes{4B} & 9 & PEQ + COMPRESSOR & BAND EMPHASIS 1 Fc~\textbf{1--2}, BAND EMPHASIS 1 Q~\textbf{1--2}, BAND EMPHASIS 1 G~\textbf{1--2}, BAND EMPHASIS 2 Fc~\textbf{3--4}, BAND EMPHASIS 2 Q~\textbf{3--4}, BAND EMPHASIS 2 G~\textbf{3--4}, THRESHOLD~\textbf{5}, RATIO~\textbf{6}, ATTACK SENSITIVITY~\textbf{7}, RELEASE SENSITIVITY~\textbf{8}, VOLUME~\textbf{9} \\
\bytes{60} & 11 & PEQ + COMPRESSOR + DISTORTION & BAND EMPHASIS 1 Fc~\textbf{1--2}, BAND EMPHASIS 1 Q~\textbf{1--2}, BAND EMPHASIS 1 G~\textbf{1--2}, BAND EMPHASIS 2 Fc~\textbf{3--4}, BAND EMPHASIS 2 Q~\textbf{3--4}, BAND EMPHASIS 2 G~\textbf{3--4}, THRESHOLD~\textbf{5}, RATIO~\textbf{6}, ATTACK SENSITIVITY~\textbf{7}, RELEASE SENSITIVITY~\textbf{8}, DRIVE~\textbf{9}, ADJUST~\textbf{10}, VOLUME~\textbf{11} \\
\bytes{61} & 9 & PEQ + COMPRESSOR + OVERDRIVE & BAND EMPHASIS Fc~\textbf{1--2}, BAND EMPHASIS Q~\textbf{1--2}, BAND EMPHASIS G~\textbf{1--2}, THRESHOLD~\textbf{3}, RATIO~\textbf{4}, ATTACK SENSITIVITY~\textbf{5}, RELEASE SENSITIVITY~\textbf{6}, DRIVE~\textbf{7}, ADJUST~\textbf{8}, VOLUME~\textbf{9} \\
\bytes{62} & 14 & PEQ + DISTORTION + DELAY & BAND EMPHASIS 1 Fc~\textbf{1--2}, BAND EMPHASIS 1 Q~\textbf{1--2}, BAND EMPHASIS 1 G~\textbf{1--2}, BAND EMPHASIS 2 Fc~\textbf{3--4}, BAND EMPHASIS 2 Q~\textbf{3--4}, BAND EMPHASIS 2 G~\textbf{3--4}, DRIVE~\textbf{5}, ADJUST~\textbf{6}, DELAY DRY/WET~\textbf{7}, DELAY L~\textbf{8--9}, DELAY R~\textbf{10--11}, FEEDBACK L~\textbf{12}, FEEDBACK R~\textbf{13}, VOLUME~\textbf{14} \\
\bytes{63} & 14 & PEQ + OVERDRIVE + DELAY & BAND EMPHASIS 1 Fc~\textbf{1--2}, BAND EMPHASIS 1 Q~\textbf{1--2}, BAND EMPHASIS 1 G~\textbf{1--2}, BAND EMPHASIS 2 Fc~\textbf{3--4}, BAND EMPHASIS 2 Q~\textbf{3--4}, BAND EMPHASIS 2 G~\textbf{3--4}, DRIVE~\textbf{5}, ADJUST~\textbf{6}, DELAY DRY/WET~\textbf{7}, DELAY L~\textbf{8--9}, DELAY R~\textbf{10--11}, FEEDBACK L~\textbf{12}, FEEDBACK R~\textbf{13}, VOLUME~\textbf{14} \\
\bottomrule
\end{longtable}
}


